% Created 2026-08-29 Sat 03:04
% Intended LaTeX compiler: pdflatex
\documentclass[11pt,twoside,a4paper]{book}
\usepackage[utf8]{inputenc}
\usepackage[T1]{fontenc}
\usepackage{graphicx}
\usepackage{longtable}
\usepackage{wrapfig}
\usepackage{rotating}
\usepackage[normalem]{ulem}
\usepackage{amsmath}
\usepackage{amssymb}
\usepackage{capt-of}
\usepackage{hyperref}
\usepackage[utf8]{inputenc}
\usepackage[T1]{fontenc}
\usepackage[a4paper,margin=1in]{geometry}
\usepackage{graphicx}
\usepackage{xcolor}
\usepackage{listings}
\usepackage{tcolorbox}
\tcbuselibrary{skins,breakable}
\usepackage{makeidx}
\makeindex
\hypersetup{colorlinks=true,linkcolor=mlmblue,citecolor=mlmblue,urlcolor=mlmblue}
\usepackage{longtable}
\usepackage{booktabs}
\bibliographystyle{plain}
\definecolor{mlmblue}{HTML}{1B4F72}
\definecolor{mlmgreen}{HTML}{1E8449}
\definecolor{mlmorange}{HTML}{B9770E}
\definecolor{kwcolor}{HTML}{7D3C98}
\definecolor{softkwcolor}{HTML}{1F618D}
\definecolor{querykwcolor}{HTML}{117864}
\definecolor{commentcolor}{HTML}{6C7A89}
\definecolor{stringcolor}{HTML}{922B21}
\definecolor{bgcode}{HTML}{F7F7F7}
\lstdefinelanguage{MLMUSE}{%
morekeywords=[1]{model,multi_model,abstract,associationClass,associationclass,%
attributes,operations,constraints,between,enum,statemachines,psm,states,%
transitions,create,context,inv,existential,pre,post,end,%
let,in,if,then,else,endif,not,and,or,xor,implies,div,true,false,null,%
Undefined,oclUndefined,oclEmpty,Set,Sequence,Bag,OrderedSet,Collection,Tuple,%
iterate,allInstances,byUseId,oclAsType,oclIsKindOf,oclIsTypeOf,selectByKind,%
selectByType,oclIsInState,oclInState,ordered,subsets,redefines,%
new,destroy,insert,into,delete,from,for,do,while,begin,declare},%
morekeywords=[2]{union,association,role,roles,composition,aggregation,dataType,%
class,signal,derived,derive,init,qualifier,MLM,mediator,clabject,assoclink},%
morekeywords=[3]{select,reject,collect,exists,forAll,isUnique,sortedBy,one,any,%
collectNested},%
sensitive=true,%
morecomment=[l]{//},%
morecomment=[l]{--},%
morecomment=[s]{/*}{*/},%
morestring=[b]',%
morestring=[b]"%
}
\lstdefinelanguage{CatMLM}{%
morekeywords=[1]{model,multi_model,abstract,associationClass,associationclass,%
attributes,operations,constraints,between,enum,statemachines,psm,states,%
transitions,create,context,inv,existential,pre,post,end,%
let,in,if,then,else,endif,not,and,or,xor,implies,div,true,false,null,%
Undefined,oclUndefined,oclEmpty,Set,Sequence,Bag,OrderedSet,Collection,Tuple,%
iterate,allInstances,byUseId,oclAsType,oclIsKindOf,oclIsTypeOf,selectByKind,%
selectByType,oclIsInState,oclInState,ordered,subsets,redefines,%
new,destroy,insert,into,delete,from,for,do,while,begin,declare},%
morekeywords=[2]{union,association,role,roles,composition,aggregation,dataType,%
class,signal,derived,derive,init,qualifier,MLM,mediator,clabject,assoclink,%
category,catAtt,catConstr,catAssociation},%
morekeywords=[3]{select,reject,collect,exists,forAll,isUnique,sortedBy,one,any,%
collectNested},%
sensitive=true,%
morecomment=[l]{//},%
morecomment=[l]{--},%
morecomment=[s]{/*}{*/},%
morestring=[b]',%
morestring=[b]"%
}
\lstdefinelanguage{ShellSession}{%
morecomment=[l]{\#},%
morestring=[b]",%
morestring=[b]',%
sensitive=true%
}
\lstdefinelanguage{EBNF}{%
morecomment=[s]{(*}{*)},%
sensitive=true%
}
\lstset{%
basicstyle=\ttfamily\small,%
keywordstyle=[1]\color{kwcolor}\bfseries,%
keywordstyle=[2]\color{softkwcolor}\bfseries,%
keywordstyle=[3]\color{querykwcolor},%
commentstyle=\color{commentcolor}\itshape,%
stringstyle=\color{stringcolor},%
backgroundcolor=\color{bgcode},%
numbers=left,numberstyle=\tiny\color{commentcolor},numbersep=8pt,%
frame=single,rulecolor=\color{black!20},%
breaklines=true,%
showstringspaces=false,%
captionpos=b,%
extendedchars=true,%
literate={—}{{---}}1,%
language=MLMUSE%
}
\newtcolorbox{manualnote}{colback=mlmblue!5,colframe=mlmblue,fonttitle=\bfseries,title=Note,breakable}
\newtcolorbox{manualwarning}{colback=mlmorange!8,colframe=mlmorange,fonttitle=\bfseries,title=Warning,breakable}
\author{MLM-USE / CatUSE Project}
\date{\today}
\title{The Multilevel USE Extensions}
\hypersetup{
 pdfauthor={MLM-USE / CatUSE Project},
 pdftitle={The Multilevel USE Extensions},
 pdfkeywords={},
 pdfsubject={},
 pdfcreator={Emacs 29.3 (Org mode 9.6.15)}, 
 pdflang={English}}
\begin{document}

\maketitle
\tableofcontents

\frontmatter
\mainmatter

\part{Getting Started}
\label{sec:org8779709}
\chapter{Introduction}
\label{sec:orgc29df65}
This manual documents a research fork of USE that adds \textbf{multi-level
modeling} (MLM-USE) and a terser front-end syntax for it (CatUSE, also
called CatMLM). It is written for someone who wants to actually write,
check, and reason about models in this system --- not just read about
the theory behind it.

\section{What Is USE?}
\label{sec:orgb19dbad}
\index{USE}

USE (the ``UML-based Specification Environment'') is a tool for
specifying and validating models written in a subset of UML together
with the Object Constraint Language (OCL). A USE specification is a
plain text file describing classes, associations, and OCL invariants;
USE compiles it, lets you populate a running \emph{system state} with
objects and links (by hand, via scripts, or interactively), and
continuously checks that state against the model's constraints.
Originally developed at the University of Bremen, it is both a
teaching tool for UML/OCL and a research platform -- this fork is an
example of the latter.

Everything in Part \ref{sec:org9ef9f58} describes plain USE as it stands
in this fork today, independent of anything multi-level.

\section{Three Layers: USE, MLM-USE, CatUSE}
\label{sec:orgd5a3a70}
\index{MLM-USE}
\index{CatUSE!relationship to MLM-USE}

This fork (developed at Ben-Gurion University, building on the
``Mediation-Based MLM in USE'' line of work) adds two more layers on top
of plain USE, each strictly \textbf{additive}: every plain USE model is
already a valid input to the layers above it.

\begin{center}
\begin{tabular}{lp{7cm}l}
Layer & Adds & Entry keyword\\[0pt]
\hline
USE & Classes, associations, OCL constraints, SOIL & \texttt{model}\\[0pt]
MLM-USE & Levels, clabjects, mediators, powerclasses & \texttt{MLM}\\[0pt]
CatUSE & Terser syntax that desugars into MLM-USE's AST & \texttt{category}\\[0pt]
\end{tabular}
\end{center}

\textbf{MLM-USE} (Part \ref{sec:org6a6c176}) is where multi-level modeling
actually lives: a class can simultaneously be an ordinary class \textbf{and}
an instance of another class one level up (a \emph{clabject}), letting a
single model describe several ontological levels -- for example, a
product catalog level whose classes are themselves instances of a
``kind of product'' level -- without collapsing them into one flat
model or duplicating information across separate models.

\textbf{CatUSE} (Part \ref{sec:org29a1bf3}, also called CatMLM in the source) is a
front-end grammar layered on top of MLM-USE. It does not add any new
semantics of its own: a CatUSE specification is parsed and then
desugared into exactly the same abstract syntax tree that an
equivalent, more verbose MLM-USE specification would produce. Compare:

\begin{lstlisting}[language=CatMLM,label= ,caption= ,captionpos=b,numbers=none,caption={CatUSE (terse)}]
category Car < Vehicle : Product
  attributes
    price : Integer
end
\end{lstlisting}

against the MLM-USE it desugars into (a separate \texttt{clabject} statement
naming the same cross-level relationship):

\begin{lstlisting}[language=MLMUSE,label= ,caption= ,captionpos=b,numbers=none,caption={The MLM-USE it desugars to}]
class Car < Vehicle
  attributes
    price : Integer
end

clabject Car : Product end
\end{lstlisting}

CatUSE currently covers less ground than MLM-USE -- some things (see
Chapter \ref{sec:org671565c}) can only be written in full
MLM-USE. Which one to reach for is mostly a matter of how much of a
model is genuinely multi-level: CatUSE reads better when most classes
in a level participate in cross-level instantiation; plain MLM-USE
keeps more control when only a few do, or when a construct CatUSE
cannot express is needed.

\section{Terminology at a Glance}
\label{sec:org80f9f29}
A handful of terms recur throughout this manual. This is only a
pointer to where each is properly introduced, not a substitute for
reading that section.

\index{clabject}
\index{mediator}
\index{powerclass}
\index{level (multi-level modeling)}

\begin{description}
\item[{\textbf{Model}}] A named collection of classes, associations, and
constraints -- the unit USE compiles and checks. See Chapter
\ref{sec:org6ab86a8}.
\item[{\textbf{Level}}] One ``floor'' of a multi-level model. See Chapter
\ref{sec:orgde23e19}.
\item[{\textbf{Clabject}}] A class that is also, simultaneously, an instance of a
class at a higher level. The name blends ``class'' and ``object''. See
Chapter \ref{sec:orgde23e19}.
\item[{\textbf{Powerclass}}] The higher-level class a clabject is an instance of.
See Chapter \ref{sec:orgde23e19}.
\item[{\textbf{Mediator}}] A declaration that relates classes across two adjacent
levels. See Chapter \ref{sec:orgde23e19}.
\item[{\textbf{Category}}] CatUSE's keyword for a classifier declaration that may
also carry cross-level instantiation directly on its own header. See
Chapter \ref{sec:org189b488}.
\end{description}

\section{How to Read This Manual}
\label{sec:org40a8231}
Part \ref{sec:org8779709} gets the system installed, built, and
running, and ends with a short guided walkthrough (Chapter
\ref{sec:orgadeff9a}). Part \ref{sec:org9ef9f58} documents plain
USE in full. Parts \ref{sec:org6a6c176} and \ref{sec:org29a1bf3} cover MLM-USE
and CatUSE respectively, each in the same shape: concepts, syntax,
worked examples. Part \ref{sec:org24bb1c9} is for anyone modifying this
fork's own implementation, not just writing models in it. Part
\ref{sec:orgec8e296} covers where USE and this fork came from. The
appendices are lookup references (grammars, keywords, shell commands),
not narrative -- go there when you need to check a detail, not to read
start to end.

Code listings are syntax-highlighted: hard (fully reserved) keywords
in one color, soft (contextual) keywords in another -- see Chapter
\ref{sec:orgfa8ba54} for what that distinction means and why it
matters. Boxed notes call out things worth pausing on:

\begin{manualnote}
This manual assumes no prior familiarity with USE. If you already know
plain USE and only want the multi-level material, Part
\ref{sec:org8779709} can be skimmed and Part \ref{sec:org9ef9f58}
skipped entirely.
\end{manualnote}
\chapter{Installing and Building}
\label{sec:org73dfb48}
\section{Prerequisites}
\label{sec:orgdc154d7}
\index{JDK}
\index{Maven}

You need:

\begin{itemize}
\item A JDK for Java 21 or later (the build sets both
\texttt{maven.compiler.source} and \texttt{maven.compiler.target} to \texttt{21}; older
JDKs cannot compile this codebase).
\item Apache Maven (any reasonably current 3.x release).
\end{itemize}

No other tools are required to build and run the system itself.
Building this manual is a separate matter, covered by its own
\texttt{Makefile} in \texttt{new-documentation/}, not by the project's Maven build.

\section{Getting the Source}
\label{sec:org0fb9f19}
\begin{lstlisting}[language=ShellSession,label= ,caption= ,captionpos=b,numbers=none]
git clone https://github.com/sturmarn/useBGU.git
cd useBGU
git checkout MLM-USE
\end{lstlisting}

All multi-level and CatUSE work described in this manual lives on the
\texttt{MLM-USE} branch of this fork, not on upstream USE.

\section{Module Layout}
\label{sec:org1de4a7e}
The build is a single Maven reactor with three modules:

\begin{center}
\begin{tabular}{lp{10cm}}
Module & Contains\\[0pt]
\hline
\texttt{use-core} & Grammars, parser, AST, the M-model (classes/associations/etc.), the compiler and system-state engine\\[0pt]
\texttt{use-gui} & The Swing GUI, the interactive shell, and the \texttt{Main} entry point\\[0pt]
\texttt{use-assembly} & Packages the other two, plus scripts and bundled examples, into a runnable distribution\\[0pt]
\end{tabular}
\end{center}

Chapters \ref{sec:orgc736412} onward go into this layout in much more
detail; this chapter only covers enough to build and run the system.

\section{Building with Maven}
\label{sec:org2671558}
From the repository root:

\begin{lstlisting}[language=ShellSession,label= ,caption= ,captionpos=b,numbers=none]
mvn clean package
\end{lstlisting}

Maven's reactor builds \texttt{use-core}, then \texttt{use-gui}, then
\texttt{use-assembly}, in that order, in a single invocation -- there is no
separate per-module build step to run by hand.

\begin{manualnote}
\texttt{mvn package} runs the unit test suite (via Surefire) as part of the
build; a failing test fails the build. To also run the slower
integration tests (Failsafe, classes named \texttt{*IT}, e.g. GUI smoke
tests), use \texttt{mvn verify} instead -- this is what CI runs.
\end{manualnote}

\section{Build Artifacts}
\label{sec:org398e54b}
After a successful build:

\begin{itemize}
\item \texttt{use-core/target/*.jar} -- the compiler/engine as a library.
\item \texttt{use-gui/target/use-gui.jar} -- a single runnable jar with all
dependencies bundled and a manifest \texttt{Main-Class} of
\texttt{org.tzi.use.main.Main}. This is the jar Chapter
\ref{sec:org6dc86d5} runs directly.
\item \texttt{use-assembly/target/use-7.1.1.tar.gz} and \texttt{.zip} -- a full
distributable archive: a \texttt{bin/} launcher script, \texttt{lib/use-gui.jar},
bundled OCL extensions, plugin jars, and the example models shipped
with the project.
\end{itemize}

The version number in the assembly's filename tracks the project's
\texttt{pom.xml} version (\texttt{7.1.1} at the time of writing); it will differ if
that version has moved on since.
\chapter{Running the System}
\label{sec:org6dc86d5}
\section{Running the Assembled Distribution}
\label{sec:org939475c}
Unpack \texttt{use-assembly/target/use-7.1.1.tar.gz} (or the \texttt{.zip}) and run
its launcher script:

\begin{lstlisting}[language=ShellSession,label= ,caption= ,captionpos=b,numbers=none]
tar xzf use-7.1.1.tar.gz
cd use-7.1.1
bin/use
\end{lstlisting}

The script resolves its own location, points at
\texttt{lib/use-gui.jar} next to it, and forwards any arguments you give it
verbatim to that jar (\texttt{bin/use [options] [spec\_file] [cmd\_file]},
covered below).

\section{Running Directly with \texttt{java -jar}}
\label{sec:org01d8df8}
Without unpacking a distribution, the freshly built jar runs on its
own:

\begin{lstlisting}[language=ShellSession,label= ,caption= ,captionpos=b,numbers=none]
java -jar use-gui/target/use-gui.jar
\end{lstlisting}

This is the same jar \texttt{bin/use} wraps, so every option below behaves
identically either way.

\section{Double-Clicking the Jar}
\label{sec:org8db5eee}
\texttt{use-gui.jar} carries a manifest \texttt{Main-Class}, so on a system where
\texttt{.jar} files are associated with a Java launcher (typically the case
on Windows and macOS out of the box; Linux desktop environments vary
and often need that association set up manually) double-clicking it
launches the GUI directly, with no arguments -- equivalent to running
\texttt{java -jar use-gui.jar} with no \texttt{spec\_file}. Passing command-line
options or a spec file to open on startup requires an actual
command-line invocation; double-clicking cannot supply either.

\section{Command-Line Options}
\label{sec:org7f068e1}
\index{command-line options}

The full usage message (\texttt{java -jar use-gui.jar -h}):

\begin{verbatim}
usage: use [options] [spec_file] [cmd_file]

options:
  -c            compile only
  -cp           compile and print specification
  -disableCollectShorthand
                flag use of OCL shorthand notation as error
  -oclAnyCollectionsChecks:W
                (W)arn|(E)rror|(I)gnore
  -extendedTypeSystemChecks:W
                (W)arn|(E)rror|(I)gnore
  -implicitTypes  Implicit variable typing in operations
  -nogui        do not use GUI
  -noplugins    do not use plugins
  -h            print help
  -H=path       home of use installation
  -nr           suppress warnings about missing readline library
  -q            reads spec_file, executes cmd_file, and checks constraints
                exit code is 1 if constraints fail, otherwise 0
  -qv           like -q but with verbose output of constraint check
  -v            print verbose messages
  -vt           print verbose messages with time info
  -V            print version info and exit

spec_file:      file containing a USE specification
cmd_file:       script file containing commands (will be executed on startup)

diagnostics:
  -debug        print lots of messages
  -p            print stack traces on errors
  -t            testing mode (less user-centric output)
\end{verbatim}

The ones you will reach for most as a new user are \texttt{-nogui} (run
headless), \texttt{-v} (see what's happening during a slow compile), \texttt{-q} /
\texttt{-qv} (batch-check a model and get a process exit code -- useful for
scripting or CI), and \texttt{-h} / \texttt{-V}. The \texttt{diagnostics} group and
\texttt{-disableCollectShorthand}, \texttt{-oclAnyCollectionsChecks}, and
\texttt{-extendedTypeSystemChecks} are aimed more at people working on the
compiler itself (Part \ref{sec:org24bb1c9}) than at someone writing
models.

\section{The Shell Runs Underneath the GUI, Too}
\label{sec:org0818ca1}
\index{shell!relationship to the GUI}

\texttt{-nogui} does not disable the interactive shell -- it only suppresses
the window. Looking at how startup actually proceeds: if a
\texttt{spec\_file} is given, it is compiled and a system state is created
\textbf{before} the GUI or shell start; then, unconditionally, a shell thread
is started; the GUI window is created only if \texttt{-nogui} was not passed.
The GUI is a window wrapped around the same underlying engine the
shell drives -- not a separate mode. Chapter
\ref{sec:org245a2e0} covers the shell's own command
language in full.

\section{Opening a Specification: Four Entry Points}
\label{sec:orgaa6bef6}
\index{opening a specification (GUI)}

The GUI's File menu offers four distinct ``Open\ldots{}'' actions, and they
are \textbf{not} content-sniffed: each parses its input with a different
grammar entry point, and you must pick the one matching what you
actually wrote.

\begin{center}
\begin{tabular}{lp{9cm}}
Action & Parses\\[0pt]
\hline
Open specification\ldots{} & A plain, single \texttt{model} (Part \ref{sec:org9ef9f58})\\[0pt]
Open multi specification\ldots{} & Several ordinary models glued by \texttt{inter-*} sections, no levels\\[0pt]
Open multi level specification\ldots{} & An \texttt{MLM} file: levels, clabjects, mediators (Part \ref{sec:org6a6c176})\\[0pt]
Open CatUSE specification\ldots{} & A CatUSE/CatMLM file (Part \ref{sec:org29a1bf3})\\[0pt]
\end{tabular}
\end{center}

Opening a file with the wrong action produces a parse error, not a
silent misinterpretation -- if in doubt, the file's own leading
keyword (\texttt{model}, \texttt{multi\_model}, \texttt{MLM}, or \texttt{category} / \texttt{class} at
the top level) tells you which one it is.

\section{Loading a Specification from the Command Line}
\label{sec:org1483098}
The command line has no separate ``which grammar'' switch -- \texttt{spec\_file}
is compiled with the entry point matching what a \texttt{-nogui} batch run
needs, inferred the same way as the table above from the file's own
syntax. A typical scripted check:

\begin{lstlisting}[language=ShellSession,label= ,caption= ,captionpos=b,numbers=none]
java -jar use-gui.jar -nogui -q myModel.use myChecks.cmd
\end{lstlisting}

\texttt{myChecks.cmd} runs on startup against the freshly loaded model; the
process exits \texttt{1} if any constraint failed, \texttt{0} otherwise -- suitable
for a CI step that should fail the build on a broken model.

\section{A Quick Sanity Check}
\label{sec:org62e108a}
To confirm the build actually works end to end, try opening one of the
worked examples that ship in the repository's own \texttt{examples/}
directory, e.g. \texttt{examples/USE/202608282106-Vehicle.use}, with ``Open
specification\ldots{}''. Chapter \ref{sec:orgadeff9a} walks through
a small model like this one in full, from writing it to checking its
first invariant.
\chapter{A Quickstart Tutorial}
\label{sec:orgadeff9a}
This chapter is a placeholder.
\part{Modeling in Plain USE}
\label{sec:org9ef9f58}
\chapter{Class Diagrams}
\label{sec:org6ab86a8}
This chapter is a placeholder.
\chapter{OCL Constraints}
\label{sec:org0cf4a3c}
This chapter is a placeholder.
\chapter{The OCL Standard Library}
\label{sec:orgd84588d}
This chapter is a placeholder.
\chapter{System States and the Shell}
\label{sec:org245a2e0}
This chapter is a placeholder.
\chapter{Validating Specifications}
\label{sec:org6928126}
This chapter is a placeholder.
\chapter{The Graphical User Interface}
\label{sec:org1b0f916}
This chapter is a placeholder.
\chapter{Model Unit Tests}
\label{sec:org4f8ed77}
This chapter is a placeholder.
\part{Multi-Level Modeling with MLM-USE}
\label{sec:org6a6c176}
\chapter{Concepts}
\label{sec:orgde23e19}
This chapter is a placeholder.
\chapter{MLM-USE Syntax in Detail}
\label{sec:orgec4571f}
This chapter is a placeholder.
\chapter{Multi-Level Shell Commands}
\label{sec:org38c2e81}
This chapter is a placeholder.
\chapter{Worked Examples}
\label{sec:org2f12cb5}
This chapter is a placeholder.
\chapter{Validating Multi-Level Models}
\label{sec:org3fbf6cb}
\index{validation!multi-level models}

This chapter is an early draft: the underlying design question --- how
\textbf{should} a multi-level model be validated, and how much of that can
reasonably be retrofitted onto tools that were never built with levels
in mind --- is still open. What follows is accurate as of this
writing, but expect it to be revised as that design settles.

Chapter \ref{sec:org6928126} covers validation for plain
USE. Everything there --- structural checking, invariant evaluation,
snapshot generation, bounded model finding --- exists as a real,
usable tool. This chapter is about what happens to each of those tools
once levels, clabjects, and inter-level constructs enter the picture.

\section{What ``Validating a Model'' Means in USE}
\label{sec:org0ba3e68}
USE actually ships several distinct things under the umbrella of
``validation,'' and it matters which one is meant:

\begin{itemize}
\item \textbf{Structural checking} --- does the model itself parse and
type-check (multiplicities, navigability, redefinition rules, and so
on)? This is what the compiler already does; nothing here is
specific to runtime state at all.
\item \textbf{Invariant evaluation against a system state} --- given a concrete,
already-populated set of objects and links, does every invariant
hold? This is ordinary OCL evaluation, nothing more.
\item \textbf{Snapshot generation} (ASSL, the \texttt{generate} shell commands) ---
backtracking search that actually \textbf{builds} a concrete population of
objects and links satisfying a set of constraints, guided by a
user-written procedure (\texttt{Create}, \texttt{CreateN}, \texttt{Insert}, \texttt{Try}, \texttt{Any},
\ldots{}). This is the closest thing USE has to ``show me a model of my
constraints.''
\item \textbf{Bounded model finding} --- a SAT/CSP-based search over a
finite-size universe, bundled with this fork as a plugin (see
\ref{sec:org668373d} in the next chapter). Conceptually similar to
Alloy's Kodkod-based analyzer.
\end{itemize}

\begin{manualnote}
\texttt{checkPartialSatisfiability()} / \texttt{checkWellDefinednessState()} --- the
two methods behind well-definedness checking --- are \textbf{not} a solver.
Each call constructs exactly one, arbitrary, concrete witness state
(every attribute defaulted to \texttt{UndefinedValue}) and evaluates the
invariant against \textbf{that single state}. It is not an existential search
over all possible states, and a \texttt{false} result does not mean ``no
witness exists'' --- only that this particular default-constructed one
doesn't satisfy it. It is easy to over-read what this check actually
proves; it proves less than ``satisfiable'' or ``unsatisfiable'' in the
logical sense.
\end{manualnote}

\section{The Multi-Level Blind Spot}
\label{sec:org18aabaf}
\index{validation!architectural blind spot}

None of the four tools above were designed with multi-level constructs
in mind, and this is not a minor omission:

\begin{itemize}
\item A clabject's \textbf{cancellation} (a removed attribute, a renamed
attribute, a removed role, a removed constraint --- see Chapter
\ref{sec:orgde23e19}) is resolved on demand, in memory, by
\texttt{MInternalClassImpl}'s own bookkeeping. A plain-USE-shaped tool that
simply walks \texttt{MClass=/=MAssociation} navigation --- which is exactly
what \texttt{MSystem} and the ASSL generator (\texttt{GGenerator}) do --- sees the
\textbf{raw}, unflattened per-level class, cancellation and all still
present.
\item \textbf{Inter-classes}, \textbf{inter-associations}, and \textbf{inter-constraints} ---
constructs owned by no single constituent model at all --- live only
in the multi-level model's own bookkeeping
(\texttt{MMultiModel.interClasses()=/=interAssociations()=/
  =interConstraints()}). A tool scoped to one constituent \texttt{MModel} at a
time never sees them.
\end{itemize}

Handed an \texttt{MMultiLevelModel} directly, \texttt{GGenerator} happily starts
generating against the \textbf{raw} per-level classes --- not an error, not a
warning, just silently wrong: cancelled attributes reappear, renamed
ones keep their original name, and inter-level constraints are never
checked at all because the generator never even looks at them. This
was confirmed directly against this fork's own \texttt{GGenerator} and
\texttt{MSystem} code, not assumed.

The bundled bounded-model-finder plugin appears, from its bytecode, to
have the identical scoping (per-\texttt{MModel}, not multi-level-aware) --- see
\ref{sec:org668373d} for the caveats on that finding, since it was
established by inspection rather than by actually running it against a
multi-level model.

\section{Two Different Ways to Think About the Problem}
\label{sec:org420d4fa}
There are two genuinely different research angles on ``validate a
model with an unbounded, or merely very large, population of linked
objects,'' worth keeping apart:

\begin{enumerate}
\item \textbf{Bounded, concrete techniques.} Snapshot generation and bounded
model finding both work by actually constructing --- or searching
over --- specific, finite object populations. They tell you
something true about the instances they build, but they say
nothing about instance counts beyond whatever bound was configured.
This is what USE actually does today, multi-level or not.
\item \textbf{Abstract, unbounded reasoning.} Shape analysis, and in particular
parametric shape analysis via three-valued logic (TVLA)
\cite{SagivRepsWilhelm:2002}, reasons about the shape of a linked
data structure --- how many objects, of what class, connected how
--- \textbf{without ever instantiating a specific population at all}. This
is the right theoretical framing for a question like ``can this
multi-level model's constraints ever be satisfied by \textbf{any} size of
population,'' but nothing in this fork attempts it; it is mentioned
here as a pointer for future work, not as a described feature.
\end{enumerate}

A related, purely theoretical point worth knowing when judging what's
achievable here at all: OCL's own expressive power depends on which
version is in play. OCL 1.4/1.5 expressions are Turing-complete;
OCL 2.0 expressions, restricted to guarantee termination, are exactly
the primitive recursive functions \cite{CengarleKnapp:2004}. This fork
was confirmed (empirically, not merely by reading the grammar) to
behave like the \textbf{unrestricted}, Turing-complete variant --- unbounded
recursion runs until a genuine \texttt{StackOverflowError} (caught and
re-thrown as a plain \texttt{RuntimeException}, not statically rejected). That
matters here because it rules out any hope of a general, exhaustive,
\textbf{static} verification procedure for arbitrary OCL constraints in this
system, multi-level or not: the underlying constraint language is, by
construction, too expressive for that. Whatever validation strategy
this fork eventually settles on for multi-level models, it will --- like
every other OCL tool of this kind --- be bounded, or heuristic, or
both, not a decision procedure.

\begin{manualwarning}
This is an open design area, not a finished feature. The next chapter
describes the one concrete step taken so far --- making the existing
bounded tools (specifically, ASSL) usable against a multi-level model
at all, by flattening it first --- and is explicit about what still
isn't wired up.
\end{manualwarning}
\chapter{The Flattening Bridge}
\label{sec:orgfee8363}
\index{validation!flattening bridge}
\index{FlatUseSystemFactory}

This chapter describes the one concrete step taken so far on the open
question raised in Chapter \ref{sec:org3fbf6cb}: making an existing
bounded validation tool --- specifically, ASSL snapshot generation ---
actually work correctly against a multi-level model, rather than
silently generating against its raw, unflattened structure. It does
not resolve the larger design question; it makes one tool usable while
that question stays open, and it is explicit below about what still
isn't done.

\section{The Idea}
\label{sec:orge254dcc}
The fix isn't to make \texttt{GGenerator} multi-level-aware --- that would be
a substantially larger undertaking than this fork's scope. It's to
hand it something it already understands correctly: an ordinary,
single-level \texttt{MModel}.

Chapter \ref{sec:orga6105c5} describes
\texttt{FlatUseTextRenderer} --- the tool that computes exactly this. Given a
compiled \texttt{MMultiLevelModel}, it renders every class's fully resolved,
post-cancellation structure (attributes, roles, invariants ---
cancellations already applied) as ordinary plain-USE text, with every
inter-class, inter-association, and inter-constraint folded in and no
multi-level syntax left anywhere in the output. Recompiling that text
with the ordinary \texttt{USECompiler} (not \texttt{USECompilerMLM}) yields a plain
\texttt{MModel} that \texttt{GGenerator} was always designed for.

The pipeline, end to end:

\begin{lstlisting}[language=MLMUSE,label= ,caption= ,captionpos=b,numbers=none]
compiled MMultiLevelModel
        |  FlatUseTextRenderer.render(...)
        v
flattened plain-USE text (a "model ... end" block, no MLM syntax)
        |  USECompiler.compileSpecification(...)
        v
plain MModel
        |  new MSystem(...)
        v
MSystem  ---  system.generator().startProcedure(...) now just works
\end{lstlisting}

\section{\texttt{FlatUseSystemFactory}}
\label{sec:org11d0f00}
\texttt{org.tzi.use.tools.catuselatex.FlatUseSystemFactory} packages the last
three steps of that pipeline as a single call:

\begin{lstlisting}[language=MLMUSE,label= ,caption= ,captionpos=b,numbers=none]
MSystem system = FlatUseSystemFactory.createFlattenedSystem(
    mlm,            // an already-compiled MMultiLevelModel
    "MyModelName",  // display name for the flattened "model" block
    err             // PrintWriter for a recompilation failure, if any
);
if (system != null) {
    system.generator().startProcedure(callString, args);
}
\end{lstlisting}

A \texttt{null} return means the flattening itself produced text that failed
to recompile --- a bug in \texttt{FlatUseTextRenderer}, not a problem with the
input model (which already compiled successfully as a multi-level
model to get this far at all).

The returned \texttt{MSystem} is a \textbf{fresh, independent} system over the
flattened model --- a different object from any \texttt{MSystem} the caller
may already have wrapping the original \texttt{mlm}. Generation results
(objects, links, invariant checks) live in the returned system, not in
whichever system \texttt{mlm} itself came from. A caller juggling both needs
to keep that distinction straight, exactly as with any two unrelated
\texttt{MSystem} instances.

\section{What's Verified}
\label{sec:org660c0c1}
\begin{itemize}
\item Every worked example under \texttt{examples/MLM-USE/} and \texttt{examples/CatMLM/}
--- including a 1000-plus-class stress-test model --- flattens and
recompiles cleanly as plain USE.
\item \texttt{FlatUseSystemFactoryTest} runs a real ASSL procedure end to end
against a flattened multi-level model via
\texttt{system.generator().startProcedure(...)}, and confirms the objects it
creates actually land in the resulting system state. Nothing in this
codebase exercised \texttt{startProcedure} against any model at all before
this test existed.
\end{itemize}

\section{The Bundled Model-Finder Plugin}
\label{sec:org668373d}
\index{validation!bounded model finding}

This fork ships a second bounded validation tool alongside ASSL:
\texttt{use-assembly/src/main/resources/plugins/ModelValidatorPlugin-5.2.0-r1.jar},
a real, bundled SAT/CSP-based model validator (built on \texttt{sat4j} and
\texttt{choco-solver}, conceptually similar to Alloy's Kodkod-based analyzer),
loaded into \texttt{use-gui} at startup via its \texttt{PluginRegistry}. It is a
pre-built binary jar shipped as a packaging resource, not a Maven
dependency of any module's own source --- worth knowing, since a plain
source or POM grep for a solver dependency will not find it.

From its bytecode, the plugin appears to be scoped the same way
\texttt{GGenerator} is: per constituent \texttt{MModel}, with no visible awareness of
clabjects, cancellation, or inter-level constructs. That finding rests
on inspection, not execution --- unlike the ASSL generator (see
\ref{sec:org660c0c1}), this plugin has not actually been run
against a multi-level model, flattened or not, so its blind spot is
inferred rather than empirically confirmed. Whether the flattening
bridge above applies to it as directly as it did to ASSL --- the
plugin has its own GUI-integrated invocation path, not a bare
\texttt{MSystem} call --- is itself an open question.

\section{What's Not Yet Wired Up}
\label{sec:orge7bf090}
\begin{manualwarning}
Two things deliberately stop here, both because they raise design
questions of their own rather than being simple wiring:

\begin{itemize}
\item \textbf{The interactive Shell.} \texttt{cmdGenStartProcedure} (the shell command
behind \texttt{gen startprocedure}) still operates directly on
\texttt{fSession.system()} --- whatever the user last \texttt{open}'d, multi-level
or not --- with no flattening step at all. Swapping in a flattened
system just for that one call raises an immediate follow-on question:
which system do the \textbf{later} generation commands (\texttt{gen result},
\texttt{gen invariant}, and friends) then operate on? The result lives in
the flattened system, not the original one, so every one of those
commands would need to know to keep using the flattened system
instead of silently looking at the wrong (and un-generated-into)
original. That's a real piece of session-state design, not yet
attempted.
\item \textbf{The bundled bounded model finder plugin} (see
\ref{sec:org668373d}) has not been run through this same bridge at
all. Its blind spot was established by bytecode inspection, not by
actually exercising it against a multi-level model the way ASSL now
has been.
\end{itemize}
\end{manualwarning}
\part{CatUSE and CatMLM}
\label{sec:org29a1bf3}
\chapter{Motivation}
\label{sec:org9b145e9}
This chapter is a placeholder.
\chapter{CatMLM Syntax in Detail}
\label{sec:org189b488}
This chapter is a placeholder.
\chapter{Known Expressiveness Gaps}
\label{sec:org671565c}
This chapter documents constructs that are deliberately unsupported, or
supported only in a restricted form, in MLM-USE and CatMLM -- as opposed
to bugs, which belong in the issue tracker, not here. Each entry says
what's restricted, why, and what to write instead. The full empirical
investigation behind each entry (reproducers, root causes in the Java
source, considered alternatives) lives in \texttt{known-issues.org} at the repo
root; this chapter gives the user-facing summary.

\section{Association classes are per-model, not multi-level}
\label{sec:org15db427}

\subsection{The restriction}
\label{sec:orge22e22d}

An \texttt{associationclass} declared directly in \texttt{inter-classes} is a full
citizen of a multi-level model. The same declaration inside an ordinary
\texttt{model ... end} block is rejected at compile time with a clear error
naming the class and model.

\subsection{Why}
\label{sec:orgd82ac31}

An association class is simultaneously a class and an association. Every
other multi-level construct keeps those two facets on entirely separate
tracks -- a clabject edge refines a class one level down, an assoclink
edge refines an association one level down, and nothing in the framework
lets one edge do both to the same declaration at once. Giving a per-model
association class real multi-level behavior would mean deciding what it
means for a single declaration to be refined as a class \textbf{and} as an
association simultaneously, including how the two refinements stay
consistent with each other -- semantics that don't exist yet in the
theory this tool implements. That is a research question, not an
implementation gap, so it is intentionally out of scope.

\subsection{What to write instead}
\label{sec:orgd01f9c8}

\begin{itemize}
\item If the association class is inherently a bridge between two different
models (or an inter-class and a model class), declare it directly in
\texttt{inter-classes} -- this already works. See
\texttt{examples/MLM-USE/202608282106-InterAssociationClassBridge.use}.
\item If it's local to one model, reify the edge as an ordinary class plus
two ordinary associations, one to each endpoint. This gives up nothing
practical: the reified edge is still an object, so population-bound
invariants (\texttt{allInstances()->size()}) apply to it exactly as they would
to a real association class. (\texttt{examples/} carries no worked file for
this pattern specifically -- it's reserved for constructs that work
without qualification, and this one exists only because of the
restriction above.)
\end{itemize}

\section{CatMLM's \texttt{instanceOf} only ever targets a level's declared parent}
\label{sec:orgfab7917}

\subsection{The restriction}
\label{sec:orgf3b0a55}

A CatMLM level is declared as \texttt{model ID1 < ID2}, and every clabject in
that level instantiates a powerclass from exactly one level: \texttt{ID2}, the
one named in that header. There is no syntax for a clabject to name a
powerclass from any other level.

\subsection{Why this isn't a gap in practice}
\label{sec:org761e5a1}

Because instanceOf is always resolved against exactly the declared
parent's class table, reusing a class name at another level -- even at
every level of a chain -- is never ambiguous: there is only ever one
table consulted, so a same-named class anywhere else is simply never
looked at. (Same-level subclassing, \texttt{<}, has the identical property for
the unrelated reason that it's resolved purely within the declaring
level's own table.) This is locked in as a regression test,
\texttt{CatMLMLevelNameDuplicationTest}, which reuses one class name at every
level of a three-level chain -- including a same-level subclass sharing
the name too -- and checks every instantiation still resolves to exactly
its own declared parent.

Note that this is a restriction of CatMLM's front-end syntax specifically,
not of MLM-USE's underlying multi-level model: a plain MLM-USE mediator
names its parent model by name (\texttt{mediator Child < Parent}), with no
adjacency requirement at all, so a lower-level class instantiating a
class from a non-adjacent ancestor level is already expressible in plain
MLM-USE today. CatMLM simply has no surface syntax to ask for that yet.

\subsection{What to write instead}
\label{sec:org6777e53}

If a class genuinely needs to instantiate a powerclass from further up
than its immediate CatMLM parent, drop to plain MLM-USE for that level
and name the target model explicitly in the \texttt{mediator ... < ...}
declaration.

\section{Other CatMLM gaps (unchanged from earlier passes)}
\label{sec:orgfc2da1e}

CatMLM's \texttt{catUSE\_model} has no \texttt{operations} section, no \texttt{inter-*}
sections (constraints/associations/classes), and no per-clabject
attribute rename -- cancellation is only via \texttt{catAtt=/=catConstr} tags on
the powerclass's own declaration, applying to every instantiator, not
selectively. Drop to plain MLM-USE for any of these.
\chapter{Worked Examples}
\label{sec:org122036d}
This chapter is a placeholder.
\chapter{The CatUSE-to-\LaTeX{} Study Tool}
\label{sec:orgbc5b501}
This chapter is a placeholder.
\part{Extending the Implementation}
\label{sec:org24bb1c9}
\chapter{Codebase Tour}
\label{sec:orgc736412}
This chapter is a placeholder.
\chapter{The Grammar and Parser}
\label{sec:org2428bfd}
This chapter is a placeholder.
\chapter{AST and Compilation Pipeline}
\label{sec:org2f8cca8}
This chapter is a placeholder.
\chapter{The Core M-Model}
\label{sec:org2126ce5}
This chapter is a placeholder.
\chapter{The Multi-Model and Multi-Level-Model Classes}
\label{sec:orgf2fe246}
This chapter is a placeholder.
\chapter{Flattening and Rendering}
\label{sec:orga6105c5}
This chapter is a placeholder.
\chapter{Testing the Compiler}
\label{sec:org81ba744}
This chapter is a placeholder.
\chapter{Releasing a New Version and Writing Plugins}
\label{sec:orgd70fe7b}
This chapter is a placeholder.
\part{History}
\label{sec:orgec8e296}
\chapter{History of USE}
\label{sec:org0e3142d}
This chapter is a placeholder.
\chapter{History of the MLM-USE/CatUSE Fork}
\label{sec:org12fe097}
This chapter is a placeholder.

\appendix

\part{Appendices}
\label{sec:org781bb22}
\chapter{The USE Grammar}
\label{sec:org6c9cadc}
The base USE language grammar (classes, associations, OCL constraints,
SOIL statements), with no multi-level-modeling and no CatMLM extensions.
Reproduced verbatim from \texttt{USE.ebnf} at the repository root, which is
generated/maintained directly from the actual grammar sources
(\texttt{USEBase.gpart}, \texttt{OCLBase.gpart}, \texttt{SoilBase.gpart},
\texttt{OCLLexerRules.gpart}) rather than transcribed by hand, so this appendix
can never drift from the real, compiled grammar.

\begin{lstlisting}[language=EBNF,label= ,caption= ,captionpos=b,numbers=none]
================================================================================
USE Language Grammar — Context-Free Grammar (EBNF)
================================================================================
This is the BASE layer: the classic, single-level USE language (classes,
associations, OCL constraints, SOIL statements) as it stands in this fork
today -- with NO multi-level-modeling and NO CatMLM extensions. It is one
of three layered EBNF files describing this fork's composed grammar:

    USE.ebnf       (this file)  -- the base language, standalone
    MLM-USE.ebnf                -- the DELTA that BGU's multi-level-modeling
                                    extension (MLM, clabject, mediator,
                                    assoclink) adds on top of this file
    CatMLM.ebnf                 -- the DELTA that the CatMLM/"CatUSE" front-
                                    end syntax sugar adds on top of MLM-USE

Each delta file only lists rules that are new or changed at that layer;
anything not mentioned there is reused unchanged from the layer(s) below.

Source: derived from use-core/src/main/resources/grammars/base/USEBase.gpart
+ OCLBase.gpart + SoilBase.gpart + OCLLexerRules.gpart, with the
multi-level-modeling rules (multi_level_model, mediator, clabject,
assoclink, and the "roles" contextual keyword) excluded -- those are
documented in MLM-USE.ebnf instead. See that file's header for why they
live in this same physical .gpart file despite being a separate layer
conceptually.

This file is a *pure* CFG: all embedded Java semantic actions, AST-builder
calls, @init blocks and ANTLR "returns [...]" clauses have been stripped
out, leaving only the production rules. Notation:

    ::=       "is defined as"
    |         alternation
    [ x ]     x is optional
    { x }     x may repeat zero or more times
    { x }+    x may repeat one or more times
    'x'       literal terminal text
    UPPERCASE reference to a lexical token (defined in the LEXICAL GRAMMAR
              section at the end)
    (* ... *) comment

Entry points (which rule the parser starts from depends on what kind of
spec is being parsed):
    model                -- a single classic USE class-diagram specification
    multi_model          -- a file containing several 'model' blocks plus
                            inter-model associations/constraints (this is
                            a base-language feature: no clabjects, no
                            mediators, no levels -- just several ordinary
                            models glued by inter-* declarations)
    expressionOnly        -- parsing a bare OCL expression (e.g. in the shell)
    typeOnly              -- parsing a bare type
    statOnly              -- parsing a bare SOIL statement

Note: 'multi_level_model' (the "MLM ..." entry point) is NOT part of this
file -- it's MLM-USE.ebnf's entry point, defined on top of multi_model_core
below (see that file).

================================================================================
SECTION 1 — CORE MODEL GRAMMAR
================================================================================

(* ---- Multi-model (several ordinary models + inter-model glue) ---- *)

multi_model ::= multi_model_core EOF

multi_model_core ::= [ 'multi_model' IDENT ]
                      internal_model { internal_model }
                      [ 'inter-enums'        { enumTypeDefinition } ]
                      [ 'inter-classes'      { interGeneralClassifierDefinition } ]
                      [ 'inter-associations' { interAssociationDefinition } ]
                      [ 'inter-constraints'  { interInvariant | interPrePost } ]

(* multi_level_model (MLM-USE.ebnf) wraps this same rule, unchanged, with
   a leading "MLM" IDENT and a trailing sequence of mediators. *)

interGeneralClassifierDefinition ::= annotationSet [ 'abstract' ]
        ( classDefinition
        | interAssociationClassDefinition
        | signalDefinition )

interAssociationClassDefinition ::= ('associationClass' | 'associationclass') IDENT
        [ '<' idList ]
        [ 'between' interAssociationEnd { interAssociationEnd }+ ]
        [ 'attributes'   { attributeDefinition } ]
        [ 'operations'   { operationDefinition } ]
        [ 'constraints'  { invariantClause } ]
        [ 'statemachines' { stateMachine } ]
        [ 'aggregation' | 'composition' ]
        'end'

interAssociationDefinition ::= annotationSet
        ('association' | 'aggregation' | 'composition') IDENT
        'between' interAssociationEnd { interAssociationEnd }+
        'end'

interAssociationEnd ::= annotationSet
        ( (IDENT '@' IDENT) | IDENT ) '[' multiplicity ']'
        [ 'role' IDENT ]
        { 'ordered'
        | 'subsets' IDENT
        | 'union'
        | 'redefines' IDENT
        | ('derived' | 'derive') [ '(' elemVarsDeclaration ')' ] '=' expression
        | 'qualifier' paramList
        }
        [ ';' ]

interInvariant ::= 'context'
        [ IDENT { ',' IDENT } ':' ]
        (multiType | simpleType)
        { invariantClause }

interPrePost ::= 'context'
        ( (IDENT '::' IDENT paramList [ ':' type ])
        | (multiType '::' IDENT paramList [ ':' type ]) )
        { prePostClause }+

(* ---- Single classic USE model ---- *)

model ::= internal_model EOF

internal_model ::= annotationSet 'model' IDENT
        { generalClassifierDefinition
        | associationDefinition
        | ( 'constraints' { invariant | prePost } )
        | enumTypeDefinition
        }

enumTypeDefinition ::= annotationSet 'enum' IDENT '{' idList '}' [ ';' ]

dataTypeDefinition ::= 'dataType' IDENT
        [ '<' idList ]
        [ 'operations'  { operationDefinition } ]
        [ 'constraints' { invariantClause } ]
        'end'

generalClassifierDefinition ::= annotationSet [ 'abstract' ]
        ( classDefinition
        | dataTypeDefinition
        | associationClassDefinition
        | signalDefinition )

classDefinition ::= 'class' IDENT
        [ '<' idList ]
        [ 'attributes'    { attributeDefinition } ]
        [ 'operations'    { operationDefinition } ]
        [ 'constraints'   { invariantClause } ]
        [ 'statemachines' { stateMachine } ]
        'end'

associationClassDefinition ::= ('associationClass' | 'associationclass') IDENT
        [ '<' idList ]
        [ 'between' associationEnd { associationEnd }+ ]
        [ 'attributes'    { attributeDefinition } ]
        [ 'operations'    { operationDefinition } ]
        [ 'constraints'   { invariantClause } ]
        [ 'statemachines' { stateMachine } ]
        [ 'aggregation' | 'composition' ]
        'end'

attributeDefinition ::= annotationSet IDENT ':' type
        [ ( (('derive'|'derived') (':'|'=') expression)
          | ('init' (':'|'=') expression) ) ]
        [ ';' ]

operationDefinition ::= annotationSet IDENT paramList superParamList
        [ ':' type ]
        [ ('=' expression) | blockStat ]
        { prePostClause }
        [ ';' ]

(* superParamList: an optional second, plain parenthesized identifier list
   right after an operation's normal parameter list, used to mark which
   parameters an operation named after its own classifier (a constructor)
   forwards to a super-classifier rather than storing as its own attribute.
   Independent of multi-level modeling -- applies to any classifier. *)
superParamList ::= [ '(' [ IDENT { ',' IDENT } ] ')' ]

associationDefinition ::= annotationSet
        ('association' | 'aggregation' | 'composition') IDENT
        'between' associationEnd { associationEnd }+
        'end'

associationEnd ::= annotationSet IDENT '[' multiplicity ']'
        [ 'role' IDENT ]
        { 'ordered'
        | 'subsets' IDENT
        | 'union'
        | 'redefines' IDENT
        | ('derived' | 'derive') [ '(' elemVarsDeclaration ')' ] '=' expression
        | 'qualifier' paramList
        }
        [ ';' ]

multiplicity      ::= multiplicityRange { ',' multiplicityRange }
multiplicityRange ::= multiplicitySpec [ '..' multiplicitySpec ]
multiplicitySpec  ::= INT | '*'

invariant ::= 'context'
        [ IDENT { ',' IDENT } ':' ]
        simpleType
        { invariantClause }

invariantClause ::= annotationSet 'inv' [ IDENT ] ':' expression
                   | 'existential' 'inv' [ IDENT ] ':' expression

prePost ::= 'context' IDENT '::' IDENT paramList [ ':' type ]
            { prePostClause }+

prePostClause ::= annotationSet ('pre' | 'post') [ IDENT ] ':' expression

annotationSet    ::= { annotation }
annotation       ::= '@' IDENT '(' annotationValues ')'
annotationValues ::= [ annotationValue ] { ',' annotationValue }
annotationValue  ::= IDENT '=' NON_OCL_STRING

stateMachine ::= 'psm' IDENT
        'states'      { stateDefinition }+
        'transitions' { transitionDefinition }+
        'end'

stateDefinition ::= IDENT [ ':' IDENT ] [ '[' expression ']' ]

transitionDefinition ::= IDENT '->' IDENT
        [ '{' [ '[' expression ']' ]
              ( event | IDENT '(' [ paramList ] ')' )
              [ '[' expression ']' ]
          '}' ]

event ::= 'create'

signalDefinition ::= 'signal' IDENT
        [ '<' idList ]
        [ 'attributes'  { attributeDefinition } ]
        [ 'constraints' { invariantClause } ]
        'end'

(* ---- Contextual ("soft") keywords ----
   The following are NOT reserved words at the lexer level: each is an
   ordinary IDENT token that is only treated as a keyword when a semantic
   predicate confirms its exact spelling at that grammar position. This
   lets 'union', 'role', 'association', 'aggregation', 'composition',
   'dataType', 'class', 'signal', 'derived', 'derive', 'init', and
   'qualifier' remain usable as ordinary identifiers elsewhere (e.g. as a
   class or attribute name) without being reserved everywhere. MLM-USE.ebnf
   and CatMLM.ebnf each add their own such keywords, scoped to their own
   grammar positions only -- see those files. *)

keyUnion       ::= IDENT   (* text must equal "union" *)
keyAssociation ::= IDENT   (* text must equal "association" *)
keyRole        ::= IDENT   (* text must equal "role" *)
keyComposition ::= IDENT   (* text must equal "composition" *)
keyAggregation ::= IDENT   (* text must equal "aggregation" *)
keyDataType    ::= IDENT   (* text must equal "dataType" *)
keyClass       ::= IDENT   (* text must equal "class" *)
keySignal      ::= IDENT   (* text must equal "signal" *)
keyDerived     ::= IDENT   (* text must equal "derived" *)
keyDerive      ::= IDENT   (* text must equal "derive" *)
keyInit        ::= IDENT   (* text must equal "init" *)
keyQualifier   ::= IDENT   (* text must equal "qualifier" *)

================================================================================
SECTION 2 — OCL EXPRESSION GRAMMAR
================================================================================

expressionOnly ::= expression EOF

expression ::= { 'let' IDENT [ ':' type ] '=' expression
                 { ',' IDENT [ ':' type ] '=' expression }
                 'in' }
               conditionalImpliesExpression

paramList ::= '(' [ variableDeclaration { ',' variableDeclaration } ] ')'

idList ::= IDENT { ',' IDENT }

variableDeclaration ::= IDENT ':' type

conditionalImpliesExpression ::= conditionalOrExpression    { 'implies' conditionalOrExpression }
conditionalOrExpression      ::= conditionalXOrExpression   { 'or' conditionalXOrExpression }
conditionalXOrExpression     ::= conditionalAndExpression   { 'xor' conditionalAndExpression }
conditionalAndExpression     ::= equalityExpression         { 'and' equalityExpression }
equalityExpression           ::= relationalExpression       { ('=' | '<>') relationalExpression }
relationalExpression         ::= additiveExpression         { ('<' | '>' | '<=' | '>=') additiveExpression }
additiveExpression            ::= multiplicativeExpression   { ('+' | '-') multiplicativeExpression }
multiplicativeExpression      ::= unaryExpression             { ('*' | '/' | 'div') unaryExpression }

unaryExpression ::= ('not' | '-' | '+') unaryExpression
                   | postfixExpression

postfixExpression ::= primaryExpression { ('.' | '->') propertyCall }

primaryExpression ::= literal
                     | objectReference
                     | propertyCall
                     | '(' expression ')'
                     | ifExpression
                     | IDENT '.' 'allInstances' [ '@' 'pre' ] [ '(' ')' ]
                     | IDENT '@' IDENT '.' 'allInstances' [ '@' 'pre' ] [ '(' ')' ]
                     | IDENT '.' 'byUseId' '(' expression ')' [ '@' 'pre' ]

objectReference ::= '@' IDENT

(* propertyCall dispatches on the identifier + following token:
     - a recognized "query" name immediately followed by '(' -> queryExpression
     - 'iterate'                                              -> iterateExpression
     - 'oclIsInState' | 'oclInState'                           -> inStateExpression
     - 'oclAsType' | 'oclIsKindOf' | 'oclIsTypeOf'
       | 'selectByType' | 'selectByKind'                       -> typeExpression
     - anything else                                           -> operationExpression
   propertyCall carries the preceding expression (the navigation source, or
   null at the start of an expression) and whether it was reached via '->'. *)

propertyCall ::= queryExpression
                | iterateExpression
                | operationExpression
                | typeExpression
                | inStateExpression

queryExpression ::= IDENT '(' [ elemVarsDeclaration '|' ] expression ')'
        (* IDENT here must be one of the recognized query operator names,
           e.g. select, reject, collect, exists, forAll, isUnique, sortedBy,
           one, any, collectNested, ... *)

iterateExpression ::= 'iterate' '('
        elemVarsDeclaration ';'
        variableInitialization '|'
        expression
        ')'

operationExpression ::= IDENT
        [ '[' expression { ',' expression } ']'
          [ '[' expression { ',' expression } ']' ] ]
        [ '@' 'pre' ]
        [ '(' [ expression { ',' expression } ] ')' ]

inStateExpression ::= ('oclIsInState' | 'oclInState') '(' IDENT ')'

typeExpression ::= ('oclAsType' | 'oclIsKindOf' | 'oclIsTypeOf'
                    | 'selectByType' | 'selectByKind') '(' type ')'

elemVarsDeclaration ::= IDENT [ ':' type ] { ',' IDENT [ ':' type ] }

variableInitialization ::= IDENT ':' type '=' expression

ifExpression ::= 'if' expression 'then' expression 'else' expression 'endif'

literal ::= 'true' | 'false'
          | INT | REAL | STRING
          | '#' IDENT
          | IDENT '::' IDENT
          | multiType '::' IDENT
          | collectionLiteral
          | emptyCollectionLiteral
          | undefinedLiteral
          | tupleLiteral
          | '*'                          (* unlimited natural, e.g. in multiplicities *)

collectionLiteral ::= ('Set' | 'Sequence' | 'Bag' | 'OrderedSet')
        '{' [ collectionItem { ',' collectionItem } ] '}'

collectionItem ::= expression [ '..' expression ]

emptyCollectionLiteral ::= 'oclEmpty' '(' collectionType ')'
                          | collectionType '{' '}'

undefinedLiteral ::= 'oclUndefined' '(' type ')'
                    | 'Undefined'
                    | 'null' '(' type ')'
                    | 'null'

tupleLiteral ::= 'Tuple' '{' tupleItem { ',' tupleItem } '}'

tupleItem ::= IDENT ( (':' type '=' expression) | ((':' | '=') expression) )

type ::= simpleType | collectionType | tupleType | multiType

typeOnly ::= type EOF

simpleType ::= IDENT
multiType  ::= IDENT '@' IDENT
        (* multiType (Level@Name) is a base-language token shape reused
           heavily by MLM-USE for fully-qualified names -- see MLM-USE.ebnf. *)

collectionType ::= ('Collection' | 'Set' | 'Sequence' | 'Bag' | 'OrderedSet')
                    '(' type ')'

tupleType ::= 'Tuple' '(' tuplePart { ',' tuplePart } ')'

tuplePart ::= IDENT ':' type

================================================================================
SECTION 3 — SOIL STATEMENT GRAMMAR (object/link manipulation language)
================================================================================

statOnly ::= stat EOF

stat ::= singleStat { ';' singleStat }

singleStat ::= emptyStat
             | statStartingWithExpr
             | varAssignStat
             | objCreateStat
             | objDestroyStat
             | lnkInsStat
             | lnkDelStat
             | condExStat
             | iterStat
             | whileStat
             | blockStat

emptyStat ::= (* empty *)

(* An expression-statement, optionally followed by an attribute
   assignment on that expression (i.e. distinguishes a bare operation
   call from  "expr.attr := rvalue"). *)
statStartingWithExpr ::= inSoilExpression [ attAssignStat ]

varAssignStat ::= IDENT ':=' rValue

attAssignStat ::= '.' IDENT ':=' rValue

objCreateStat ::= 'new' simpleType
        [ '(' [ inSoilExpression ] ')' ]
        [ 'between' '(' rValListMin2WithOptionalQualifiers ')' ]

objDestroyStat ::= 'destroy' exprListMin1

lnkInsStat ::= 'insert' '(' rValListMin2WithOptionalQualifiers ')' 'into'
               ( IDENT | IDENT '@' IDENT )

rValListMin2WithOptionalQualifiers ::=
        rValue [ '{' rValList '}' ]
        ',' rValue [ '{' rValList '}' ]
        { ',' rValue [ '{' rValList '}' ] }

lnkDelStat ::= 'delete' '(' rValListMin2WithOptionalQualifiers ')' 'from' IDENT

condExStat ::= 'if' inSoilExpression 'then' statOrImplicitBlock
               [ 'else' statOrImplicitBlock ]
               'end'

iterStat ::= 'for' IDENT 'in' inSoilExpression 'do' statOrImplicitBlock 'end'

whileStat ::= 'while' inSoilExpression 'do' statOrImplicitBlock 'end'

blockStat ::= 'begin'
        [ 'declare' variableDeclaration { ',' variableDeclaration } ';' ]
        stat
        'end'

implicitBlockStat ::= 'declare' variableDeclaration { ',' variableDeclaration } ';' stat

statOrImplicitBlock ::= stat | implicitBlockStat

rValue ::= inSoilExpression | objCreateStat

rValList     ::= [ rValListMin1 ]
rValListMin1 ::= rValue { ',' rValue }
rValListMin2 ::= rValue ',' rValue { ',' rValue }

inSoilExpression ::= expression

exprList     ::= [ exprListMin1 ]
exprListMin1 ::= inSoilExpression { ',' inSoilExpression }
exprListMin2 ::= inSoilExpression ',' inSoilExpression { ',' inSoilExpression }

identList     ::= [ identListMin1 ]
identListMin1 ::= IDENT { ',' IDENT }

================================================================================
SECTION 4 — LEXICAL GRAMMAR (tokens)
================================================================================
Shared verbatim by all three layers (USE, MLM-USE, CatMLM) -- none of them
add or change a token definition.

(* Whitespace / comments -- discarded (sent to the HIDDEN channel) *)
WS          ::= ' ' | '\t' | '\f' | NEWLINE
SL_COMMENT  ::= ('//' | '--') { any char except NEWLINE }
ML_COMMENT  ::= '/*' { any char, non-greedy } '*/'
NEWLINE     ::= '\r\n' | '\r' | '\n'

(* Punctuation / operator tokens *)
ARROW         ::= '->'
AT            ::= '@'
BAR           ::= '|'
COLON         ::= ':'
COLON_COLON   ::= '::'
COLON_EQUAL   ::= ':='
COMMA         ::= ','
DOT           ::= '.'
DOTDOT        ::= '..'
EQUAL         ::= '='
GREATER       ::= '>'
GREATER_EQUAL ::= '>='
HASH          ::= '#'
LBRACE        ::= '{'
LBRACK        ::= '['
LESS          ::= '<'
LESS_EQUAL    ::= '<='
LPAREN        ::= '('
MINUS         ::= '-'
NOT_EQUAL     ::= '<>'
PLUS          ::= '+'
RBRACE        ::= '}'
RBRACK        ::= ']'
RPAREN        ::= ')'
SEMI          ::= ';'
SLASH         ::= '/'
STAR          ::= '*'

(* Numbers *)
INT             ::= digit { digit }                                    (* digit = '0'..'9' *)
REAL            ::= INT '.' INT [ ('e'|'E') ['+'|'-'] INT ]
                   | INT ('e'|'E') ['+'|'-'] INT
RANGE_OR_INT    ::= INT (retyped to INT if followed by '..', else to REAL if it forms a REAL, else stays INT)

(* Strings *)
STRING          ::= "'" { any char except "'" or '\' | ESC } "'"           (* OCL string literal *)
NON_OCL_STRING  ::= '"' { any char except '"' or '\' | ESC } '"'           (* used only in annotation values *)
ESC             ::= '\' ( 'n'|'r'|'t'|'b'|'f'|'"'|"'"|'\'
                         | 'u' HEX_DIGIT HEX_DIGIT HEX_DIGIT HEX_DIGIT
                         | '0'..'3' [ '0'..'7' [ '0'..'7' ] ]
                         | '4'..'7' [ '0'..'7' ] )
HEX_DIGIT       ::= '0'..'9' | 'A'..'F' | 'a'..'f'

(* Identifiers *)
IDENT ::= ('$' | letter | '_') { letter | '_' | digit }
        (* letter = 'a'..'z' | 'A'..'Z' *)
        (* After matching, ANTLR's literal table reclassifies IDENT tokens
           whose text matches a quoted keyword used elsewhere in the
           grammar (e.g. 'class', 'model', 'context', 'inv', 'let', ...)
           as that keyword token instead. The contextual "soft" keywords
           listed above (union, role, association, ...) are NOT handled
           this way -- they stay IDENT and are matched only via semantic
           predicate at the point of use. *)

================================================================================
NOTES
================================================================================
1. This is a *pure* CFG view of the grammar for reading/reasoning purposes.
   The real ANTLR3 grammar (use-core/target/grammars/USE.g, regenerated on
   every `mvn compile`/`mvn package` from the .gpart sources) is not
   equivalent to this file in the formal sense: it interleaves embedded
   Java actions that build the AST as parsing proceeds, uses a few
   syntactic-predicate lookaheads (e.g. `(statStartingWithExpr)=>`,
   `(COLON type EQUAL)=>`) to resolve local ambiguities that a plain CFG
   can't express, and relies on semantic predicates for the contextual
   keywords noted above. Treat this file as documentation, not as
   something you could feed straight into a parser generator.

2. This file, together with MLM-USE.ebnf and CatMLM.ebnf, reconstructs the
   full USE.g composition (USEBase.gpart + OCLBase.gpart + SoilBase.gpart +
   OCLLexerRules.gpart, layered with the multi-level and CatMLM additions)
   -- the grammar for ordinary .use / .mlm specification files, not the
   interactive shell-command grammar (a separate, sixth fragment,
   ShellCommandBase.gpart, builds that one instead).
\end{lstlisting}
\chapter{The MLM-USE Grammar}
\label{sec:org18e3b8a}
The multi-level-modeling extension, as a \textbf{delta} on top of the USE
grammar in Appendix \ref{sec:org6c9cadc}: only rules that are new or
changed at this layer are listed. Reproduced verbatim from
\texttt{MLM-USE.ebnf} at the repository root.

This grammar describes syntax only. A construct can parse cleanly here
and still be rejected during semantic analysis once it's nested inside a
multi-level model -- most notably, a \texttt{model ... end} block's classifier
list reuses the plain USE \texttt{associationclass} rule unchanged, so an
association class inside a \texttt{model} block parses, but is rejected at
compile time (it's only accepted in \texttt{inter-classes}). See Chapter
\ref{sec:org671565c} for this and the other cases where the
grammar alone doesn't tell the whole story.

\begin{lstlisting}[language=EBNF,label= ,caption= ,captionpos=b,numbers=none]
================================================================================
MLM-USE Grammar — DELTA over USE.ebnf (EBNF)
================================================================================
This file lists ONLY what the BGU multi-level-modeling extension (clabjects,
mediators, powerclasses, assoclinks) adds on top of the base language in
USE.ebnf. Anything not mentioned here -- class/association/constraint/OCL/
SOIL syntax, the "model"/"multi_model" entry points, etc. -- is exactly the
grammar in USE.ebnf, unchanged.

Notation is the same as USE.ebnf:
    ::=       "is defined as"        |    alternation
    [ x ]     x is optional          { x }    x may repeat zero or more times
    { x }+    x may repeat one or more times
    'x'       literal terminal text
    UPPERCASE reference to a lexical token (see USE.ebnf Section 4)
    (* ... *) comment

--------------------------------------------------------------------------------
Where this actually lives (important, and the reason the old merged
USE-grammar.ebnf file was misleading about it): despite the project's
directory layout suggesting a clean per-feature split --

    use/USE.gpart              -- (its own header comment claims:) "multi-
                                   level modeling additions (MLM, clabject,
                                   mediator, assoclink) + core UML model
                                   syntax"
    base/USEBase.gpart         -- "core class/association/constraint syntax"

-- use/USE.gpart in fact contains ONLY grammar boilerplate (package
declarations, lexer error-handling members) and defines NO production
rules at all. Every rule below is physically defined inside
base/USEBase.gpart, interleaved with the base-language rules from
USE.ebnf. Confirmed by diffing USEBase.gpart against its state at commit
6106270f, the last commit before MLM work began (571a84bb, "Began
implementing MLM, added parsing rules and AST classes") -- multi_level_
model, mediator, clabject, and assoclink are a clean, self-contained
insertion with no further changes anywhere else in the file since. This
file corrects that attribution: everything below is the true MLM-USE
delta, regardless of which physical .gpart file it happens to sit in.
--------------------------------------------------------------------------------

New entry point (alongside 'model' and 'multi_model' from USE.ebnf):
    multi_level_model   -- a full multi-level model file (starts with 'MLM')

================================================================================
NEW RULES
================================================================================

multi_level_model ::= 'MLM' IDENT multi_model_core { mediator } EOF
        (* multi_model_core is USE.ebnf's rule, reused completely unchanged
           -- an MLM file is exactly "MLM <name>" followed by the same
           several-models-plus-inter-* body multi_model already accepts,
           followed by zero or more mediators. "MLM" here is a soft keyword
           (keyMLM below), not a global reserved word -- see that section. *)

mediator ::= 'mediator' IDENT '<' IDENT
             { clabject | assoclink }
             'end'
        (* "mediator" is a soft keyword (keyMediator below). *)

clabject ::= 'clabject' IDENT ':' IDENT
             [ 'attributes' { (IDENT '->' IDENT) | ('~' IDENT) } ]
             [ 'roles'      { '~' IDENT } ]
             [ 'constraints' { '~' IDENT } ]
             'end'
        (* The child (IDENT before ':') instantiates the powerclass (IDENT
           after ':') from the mediator's parent model. The three optional
           sections cancel or rename inherited features: attribute renaming
           (IDENT -> IDENT) or removal (~IDENT), role removal, and
           constraint removal -- each by the powerclass's own feature name.
           "clabject" is a soft keyword (keyClabject below); "roles" is too
           (keyRoles below). *)

assoclink ::= 'assoclink' IDENT ':' IDENT
              (IDENT '->' IDENT)
              (IDENT '->' IDENT)
              'end'
        (* An assoclink between two clabjects instantiates an association
           between their two powerclasses; the two role-renaming pairs
           bind the child association's two role names to the parent
           association's two role names. "assoclink" is a soft keyword
           (keyAssoclink below). *)

================================================================================
NEW CONTEXTUAL ("soft") KEYWORDS
================================================================================
keyMLM       ::= IDENT   (* text must equal "MLM"; matched only at the start
                            of a multi_level_model *)
keyMediator  ::= IDENT   (* text must equal "mediator"; matched only at a
                            mediator block's header *)
keyClabject  ::= IDENT   (* text must equal "clabject"; matched only at a
                            clabject declaration's header *)
keyAssoclink ::= IDENT   (* text must equal "assoclink"; matched only at an
                            assoclink declaration's header *)
keyRoles     ::= IDENT   (* text must equal "roles"; matched only inside the
                            clabject rule's own "roles" section header above *)

None of these five existed as soft keywords from the start. All five were
originally written as bare ANTLR3 literals directly in these rules, which
auto-promotes each one to a keyword reserved *everywhere* in its composed
grammar -- not just inside the one multi-level-modeling section it actually
belongs to. That silently broke any plain, non-MLM USE model that happened
to use "roles", "MLM", "mediator", "clabject", or "assoclink" as an ordinary
identifier (e.g. an OCL operation parameter, or an attribute name) -- "roles"
was caught first, breaking use-gui's own ShellIT/t109.use integration
fixture; the other four were found and fixed as a direct follow-up, confirmed
empirically the same way (a plain model with an attribute named "clabject"
failed to parse with "mismatched input 'clabject' expecting 'end'" before
the fix). "MLM" and "clabject" were bare literals in CatMLM.gpart too,
independently of USEBase.gpart -- since CatMLM is a separately-composed
grammar with its own lexer, both files needed fixing; CatMLM.gpart's rules
now reuse the keyMLM/keyClabject rules defined once here, the same way they
already reused keyClass and keyAssociation from USE.ebnf. Fixed by
introducing these five soft keywords, the same pattern USE.ebnf's base
contextual keywords already use (keyRole, keyUnion, keyAssociation, ...):
confirmed via an isolated git worktree that the "roles" instance of the bug
predated any multi-level-modeling work and was purely a base-grammar defect
introduced alongside the clabject rule.

================================================================================
NOTES
================================================================================
1. Every non-terminal referenced above but not defined here or reused
   verbatim from USE.ebnf (IDENT, and the multi_model_core it wraps) comes
   from USE.ebnf. multiType (IDENT '@' IDENT, i.e. "Level@Name") is USE.ebnf's
   token shape but is where MLM-USE's fully-qualified names live at the
   semantic level -- the grammar itself doesn't need a separate production
   for it since ordinary type/simpleType references already parse it.

2. As with USE.ebnf, this is a pure-CFG view for reading/reasoning; the real
   grammar interleaves Java AST-builder actions (e.g. resolving a mediator's
   parent model by name, not by declaration order -- a bug found and fixed
   in this fork's Java layer, not visible at the grammar level at all).

3. See CatMLM.ebnf for the front-end sugar layer built on top of this one --
   CatMLM's "category"/"clabject" shorthand and cat* tags desugar, before
   any of this grammar runs, into exactly the plain MLM-USE AST that parsing
   the rules above would itself produce.
\end{lstlisting}
\chapter{The CatMLM Grammar}
\label{sec:org529f59b}
The CatMLM (``CatUSE'') front-end syntax sugar, as a \textbf{delta} on top of
MLM-USE (Appendix \ref{sec:org18e3b8a}): CatMLM desugars into
exactly the same AST that plain MLM-USE produces. Reproduced verbatim
from \texttt{CatMLM.ebnf} at the repository root.

\begin{lstlisting}[language=EBNF,label= ,caption= ,captionpos=b,numbers=none]
================================================================================
CatMLM ("CatUSE") Grammar — DELTA over MLM-USE.ebnf (EBNF)
================================================================================
This file lists ONLY what the CatMLM front-end -- the terser "category"/
"clabject" surface syntax with cat* tags, informally called "CatUSE" -- adds
on top of MLM-USE.ebnf. It is a separate, additive grammar: it changes
nothing in USEBase.gpart/USE.ebnf/MLM-USE.ebnf, and everything downstream
of parsing (MMultiLevelModel, well-definedness checking, the GUI, ...) has
no notion that CatMLM exists at all. A CatMLM parse tree is *desugared*, in
Java (USECompilerCatUSE.desugar), into the exact same ASTMultiLevelModel /
ASTMediator / ASTClabject / ASTModel / ASTClass tree that parsing the
equivalent plain MLM-USE.ebnf source would produce.

Source: use-core/src/main/resources/grammars/catmlm/CatMLM.gpart, composed
as CatMLM.gpart + USEBase.gpart + OCLBase.gpart + SoilBase.gpart +
OCLLexerRules.gpart -- a genuinely separate ANTLR grammar/parser
(CatUSEParser) from the one USE.ebnf/MLM-USE.ebnf describe, sharing only
the reused-unchanged rules noted below.

Notation is the same as USE.ebnf:
    ::=       "is defined as"        |    alternation
    [ x ]     x is optional          { x }    x may repeat zero or more times
    { x }+    x may repeat one or more times
    'x'       literal terminal text
    UPPERCASE reference to a lexical token (see USE.ebnf Section 4)
    (* ... *) comment

New entry point:
    catUSE_model   -- a full CatMLM file (also starts with 'MLM' IDENT --
                      see the dispatch note below)

Note on dispatch: catUSE_model and MLM-USE.ebnf's multi_level_model both
start with the literal sequence 'MLM' IDENT, so a bare file is NOT
self-describing -- there is no sniffing of file content to decide which
grammar to parse it with. The tool chooses: the GUI has a separate "Open
CatUSE specification..." action (MainWindow.ActionFileOpenCatUSESpec)
distinct from its plain "Open MLM specification..." action, and
USECompilerCatUSE.compileCatUSESpecification vs. the plain MLM compiler
are two different entry points at the Java level.

================================================================================
NEW RULES
================================================================================

catUSE_model ::= 'MLM' IDENT { catLevel }+ EOF
        (* "MLM" here is MLM-USE.ebnf's keyMLM soft keyword, reused
           unchanged -- not redefined in this file, see below. *)

catLevel ::= 'model' IDENT [ '<' IDENT ]
             { catClassifierDefinition
             | catAssociationDefinition
             | catClabjectDefinition
             | ( 'constraints' { catInvariant } )
             }
        (* Fuses what MLM-USE.ebnf splits across a separate 'model' block
           (in multi_model_core) and a separate 'mediator' block: one
           "model ID [< Parent]" section holds both the level's own
           classifiers/associations/constraints AND its clabjects (which
           desugar into the *parent* level's mediator). *)

catClassifierDefinition ::= ( 'category' | 'class' ) IDENT
        [ '<' catRefList ]
        [ ':' catRefList ]
        [ 'attributes'  { catAttributeDefinition } ]
        [ 'constraints' { catInvariantClause } ]
        'end'
        (* "category X ... end" and "class X ... end" desugar to the exact
           same AST class either way -- isCategory is purely a hint for a
           human reader (and, incidentally, for the study tool's PlantUML
           diagram stereotypes); it is never checked against anything and
           decides nothing about the resulting model's semantics.

           The ":" clause is Java-"implements"-style fusion of instantiation
           onto the classifier's own header, e.g.:
             category C2 < D : C3
             attributes
               seats: Integer
             end
           as an alternative to a separate catClabjectDefinition statement
           for the same name. It desugars to exactly the same ASTCatClabject
           a standalone "clabject C2 : C3 end" would produce, added to the
           same level -- so the two spellings are sugar for each other, and
           using both for the same class name is not an error: their
           reference lists union (see the NOTES section on silent union). *)

catRefList ::= IDENT
             | '(' IDENT ',' IDENT { ',' IDENT } ')'
        (* One bare identifier, or two-or-more wrapped in parentheses --
           used for every class-name reference list in CatMLM: both
           clauses of catClassifierDefinition above, and
           catClabjectDefinition's powerclass list below. The grammar
           itself has no ambiguity that requires the parentheses (":"
           already unambiguously ends a preceding "<" list on its own,
           and a standalone list never sits next to another one) -- this
           is a readability convention, applied consistently everywhere
           such a list can appear, not a parsing necessity. A single name
           is never parenthesized: "(A)" is rejected, only "A" or
           "(A, B, ...)" are valid. *)

catAttributeDefinition ::= IDENT ':' type [ ';' 'catAtt' ] [ ';' ]
        (* Reuses IDENT/COLON/type exactly like USE.ebnf's attributeDefinition;
           only the optional trailing "catAtt" tag is new. A catAtt-tagged
           attribute is automatically cancelled in every clabject that
           instantiates this (or an inheriting) category. *)

catInvariantClause ::= 'inv' [ IDENT ] ':' [ 'catConstr' ] expression
        (* The "catConstr" tag sits between the colon and the expression.
           A catConstr-tagged invariant is automatically cancelled in every
           instantiating clabject, the same way catAtt cancels attributes. *)

catInvariant ::= 'context'
        [ IDENT { ',' IDENT } ':' ]
        simpleType
        { catInvariantClause }
        (* Mirrors USE.ebnf's invariant rule exactly, just built from
           catInvariantClause instead of invariantClause. *)

catAssociationDefinition ::= ( 'catAssociation' | 'association' | 'aggregation' | 'composition' ) IDENT
        'between' associationEnd { associationEnd }+
        'end'
        (* associationEnd is MLM-USE.ebnf/USE.ebnf's rule, reused unchanged.
           A catAssociation-tagged association has its far-end role(s)
           cancelled for any clabject instantiating the near end -- this is
           what stops a clabject instantiating two powerclasses from ending
           up with two different roles that happen to share a name. *)

catClabjectDefinition ::= 'clabject'
        [ IDENT '<' IDENT ';' ]
        IDENT ':' catRefList
        'end'
        (* One declaration, N powerclasses, and an optional fused same-level
           subclass edge, all collapsing to this single rule:
             clabject C : A end
             clabject C : (A, B) end
             clabject E < D; E : B2 end
           The first optional "IDENT < IDENT ;" is a same-level superclass
           edge (like MLM-USE.ebnf's own "class X < Y"); the required
           "IDENT : catRefList" lists the one-or-more powerclasses this
           clabject instantiates (parenthesized if more than one -- see
           catRefList above), desugaring to one plain MLM-USE
           `clabject <name> : <powerclass>` per listed powerclass. No
           keyword precedes the powerclass list -- see below. The leading
           "clabject" is MLM-USE.ebnf's keyClabject soft keyword, reused
           unchanged -- not redefined in this file, see below. *)

================================================================================
NEW CONTEXTUAL ("soft") KEYWORDS
================================================================================
keyCategory       ::= IDENT   (* text must equal "category"; matched only
                                 in catClassifierDefinition, choosing between
                                 "category ID ... end" and "class ID ... end" *)
keyCatAtt         ::= IDENT   (* text must equal "catAtt" *)
keyCatConstr      ::= IDENT   (* text must equal "catConstr" *)
keyCatAssociation ::= IDENT   (* text must equal "catAssociation" *)

(keyMLM and keyClabject, used above, are MLM-USE.ebnf's soft keywords --
CatMLM.gpart reuses those rule definitions rather than declaring its own
"MLM"/"clabject" literals a second time, which would have created the same
"reserved everywhere in this composed grammar's lexer" problem independently
of USEBase.gpart. This mirrors how catClassifierDefinition already reuses
USE.ebnf's keyClass, and catAssociationDefinition reuses keyAssociation,
keyAggregation, and keyComposition.)

Note on "category" and catClabjectDefinition: an earlier version of this
grammar also accepted an optional "category" hint immediately after the
colon in catClabjectDefinition (i.e. "clabject C : category A, B end").
It was removed: the "<" vs ":" already unambiguously distinguishes
same-level generalization from cross-level instantiation, so the keyword
added no information a reader didn't already have from the symbol choice
-- and it was never actually checked against the named class's own
declaration (a clabject could "instantiate" a plain "class", not just a
"category", keyword or no keyword, both before and after this change).
keyCategory today is used exclusively in catClassifierDefinition.

================================================================================
NOTES
================================================================================
1. Every non-terminal referenced above but not defined here -- type,
   expression, simpleType, associationEnd -- is reused completely unchanged
   from USE.ebnf/MLM-USE.ebnf. catRefList is CatMLM's own rule (defined
   above): the base grammar's plain idList (a bare comma list, no
   parentheses ever) is no longer used anywhere in CatMLM.gpart at all --
   catRefList replaced its two former uses (catClassifierDefinition's "<"
   list, catClabjectDefinition's ":" list) and is also used for the new
   ":" clause on catClassifierDefinition.

2. As with the other two files, this is a pure-CFG view for reading and
   reasoning, not something to feed straight into a parser generator: the
   real grammar's semantic actions build the desugared AST as parsing
   proceeds (see USECompilerCatUSE.desugar's four translation rules), and
   ANTLR's adaptive lookahead resolves a few local ambiguities a plain CFG
   can't express -- for instance, a category or clabject literally named
   "category" (or "catAtt", "role", etc.) still parses correctly as an
   ordinary identifier wherever an IDENT is expected, backtracking off the
   soft-keyword reading when it can't also complete a valid parse.

3. Silent union, and a bug it uncovered: declaring the same class name's
   instantiation or same-level superclassifiers via more than one route --
   the fused ":"/"<" header on catClassifierDefinition, a separate
   catClabjectDefinition statement, and/or that statement's own fused
   "IDENT < IDENT ;" local-superclass prefix -- is not an error. All of
   them contribute to the same union, by design: e.g. a classifier declared
   with "category C2 < D : C3 end" and, elsewhere in the same level, a
   separate "clabject C2 : C4 end" both apply to the one class C2, which
   ends up instantiating both C3 and C4. Building this surfaced a genuine
   pre-existing bug (not new in this change) in
   ASTClassifier.addSuperClassifiers: it used to *replace*
   (fSuperClassifiers = idList) rather than accumulate, so a class whose
   same-level superclassifiers were set from two different declaration
   sites would silently keep only whichever one ran last. Fixed to
   accumulate instead -- safe for every base-grammar caller too, since none
   of them ever call it more than once per classifier.
\end{lstlisting}
\chapter{Keyword Reference}
\label{sec:orgfa8ba54}
This chapter is a placeholder.
\chapter{Shell Command Reference}
\label{sec:org1586b85}
This chapter is a placeholder.

\backmatter
\bibliography{bib/references}
\printindex
\end{document}
